
\documentclass{article}
\usepackage{colortbl}
\usepackage{makecell}
\usepackage{multirow}
\usepackage{supertabular}

\begin{document}

\newcounter{utterance}

\centering \large Interaction Transcript for game `dond', experiment `semi', episode 9 with Qwen3.6{-}35B{-}A3B{-}FP8{-}without{-}reasoning.
\vspace{24pt}

{ \footnotesize  \setcounter{utterance}{1}
\setlength{\tabcolsep}{0pt}
\begin{supertabular}{c@{$\;$}|p{.15\linewidth}@{}p{.15\linewidth}p{.15\linewidth}p{.15\linewidth}p{.15\linewidth}p{.15\linewidth}}
    \# & $\;$A & \multicolumn{4}{c}{Game Master} & $\;\:$B\\
    \hline

    \theutterance \stepcounter{utterance}  
    & & \multicolumn{4}{p{0.6\linewidth}}{
        \cellcolor[rgb]{0.9,0.9,0.9}{
            \makecell[{{p{\linewidth}}}]{
                \texttt{\tiny{[P1$\langle$GM]}}
                \texttt{Igrate igro sodelovalnega pogajanja, v kateri se morate dogovoriti s drugim igralcem o tem, kako razdeliti nabor predmetov med seboj.} \\
\\ 
\texttt{Pravila:} \\
\texttt{(a) Ti in drugi igralec dobite nabor predmetov. Vsakemu od vas je dana tudi tajna funkcija vrednosti, ki predstavlja, kako zelo cenite vsako vrsto predmeta.} \\
\texttt{(b) Izmenjujete si sporočila z drugim igralcem, da se dogovorite, kdo dobi katere predmete. Vsak lahko pošlje največ 5 sporočil ali zaključi igro prej tako, da v katerem koli trenutku daš skritega predloga.} \\
\texttt{(c) Prosimo vas, da pošljete skritega predloga, ki navaja predmete, ki jih želite, oblikovan v oglatih oglatih oklepajih, kot sledi: "{[}Predlog: <število> <ime predmeta>, <število> <ime predmeta>, <...>{]}"} \\
\texttt{(d) Če so vaši predlogi dopolnilni, tj. je dovolj predmetov za izpolnitev obeh predlogov, vsak igralec dobi rezultat na podlagi vsote vrednosti predmetov, ki jih je prejel. V nasprotnem primeru oba dobita nič točk.} \\
\texttt{(e) Vaš cilj je povečati vsoto vašega rezultata in rezultata drugega igralca.} \\
\\ 
\texttt{Začnimo.} \\
\\ 
\texttt{Niz razpoložljivih predmetov je:} \\
\\ 
\texttt{1 lonec, 3 cvetovi, 1 baterija} \\
\\ 
\texttt{Vaša tajna funkcija vrednosti je:} \\
\\ 
\texttt{lonec: 5, cvet: 0, baterija: 5} \\
\\ 
\texttt{POMEMBNO: Vaša sporočila, razen če gre za skriti predlog, se posredujejo neposredno drugemu igralcu, zato ne vključujte nobenega odgovora na pravila ali besedilo, ki napoveduje vaše sporočilo. Če želite podati skriti predlog, uporabite navedeni format. Ne uporabljajte oglatih oglatih oklepajev pri komunikaciji z drugim igralcem, saj bo to razumljeno kot vaš skritega predloga.} \\
            }
        }
    }
    & \\ \\

    \theutterance \stepcounter{utterance}  
    & \multicolumn{4}{p{0.6\linewidth}}{
        \cellcolor[rgb]{0.9,0.9,0.9}{
            \makecell[{{p{\linewidth}}}]{
                \texttt{\tiny{[P1$\rangle$GM]}}
                \texttt{Pozdravljen! Predlagam, da se dogovorimo za razdelitev, ki bo obema igralcema prinesla čim večjo skupno vrednost. Glede na to, da imam jaz visoko vrednost za lonec in baterijo, predlagam naslednjo porazdelitev:} \\
\\ 
\texttt{Jaz bi rad lonec in baterijo, ti pa lahko vzameš vse tri cvetove.} \\
\\ 
\texttt{Prosim, sporoči mi svoj predlog ali potrdi to razdelitev.} \\
            }
        }
    }
    & & \\ \\

    \theutterance \stepcounter{utterance}  
    & & & \multicolumn{2}{p{0.3\linewidth}}{
        \cellcolor[rgb]{0.9,0.9,0.9}{
            \makecell[{{p{\linewidth}}}]{
                \texttt{\tiny{[GM$|$GM]}}
                \texttt{continue} \\
            }
        }
    }
    & & \\ \\

    \theutterance \stepcounter{utterance}  
    & & & \multicolumn{2}{p{0.3\linewidth}}{
        \cellcolor[rgb]{0.9,0.9,0.9}{
            \makecell[{{p{\linewidth}}}]{
                \texttt{\tiny{[GM$|$GM]}}
                \texttt{Pozdravljen! Predlagam, da se dogovorimo za razdelitev, ki bo obema igralcema prinesla čim večjo skupno vrednost. Glede na to, da imam jaz visoko vrednost za lonec in baterijo, predlagam naslednjo porazdelitev:} \\
\\ 
\texttt{Jaz bi rad lonec in baterijo, ti pa lahko vzameš vse tri cvetove.} \\
\\ 
\texttt{Prosim, sporoči mi svoj predlog ali potrdi to razdelitev.} \\
            }
        }
    }
    & & \\ \\

    \theutterance \stepcounter{utterance}  
    & & \multicolumn{4}{p{0.6\linewidth}}{
        \cellcolor[rgb]{0.9,0.9,0.9}{
            \makecell[{{p{\linewidth}}}]{
                \texttt{\tiny{[GM$\rangle$P2]}}
                \texttt{Igrate igro sodelovalnega pogajanja, v kateri se morate dogovoriti s drugim igralcem o tem, kako razdeliti nabor predmetov med seboj.} \\
\\ 
\texttt{Pravila:} \\
\texttt{(a) Ti in drugi igralec dobite nabor predmetov. Vsakemu od vas je dana tudi tajna funkcija vrednosti, ki predstavlja, kako zelo cenite vsako vrsto predmeta.} \\
\texttt{(b) Izmenjujete si sporočila z drugim igralcem, da se dogovorite, kdo dobi katere predmete. Vsak lahko pošlje največ 5 sporočil ali zaključi igro prej tako, da v katerem koli trenutku daš skritega predloga.} \\
\texttt{(c) Prosimo vas, da pošljete skritega predloga, ki navaja predmete, ki jih želite, oblikovan v oglatih oglatih oklepajih, kot sledi: "{[}Predlog: <število> <ime predmeta>, <število> <ime predmeta>, <...>{]}"} \\
\texttt{(d) Če so vaši predlogi dopolnilni, tj. je dovolj predmetov za izpolnitev obeh predlogov, vsak igralec dobi rezultat na podlagi vsote vrednosti predmetov, ki jih je prejel. V nasprotnem primeru oba dobita nič točk.} \\
\texttt{(e) Vaš cilj je povečati vsoto vašega rezultata in rezultata drugega igralca.} \\
\\ 
\texttt{Začnimo.} \\
\\ 
\texttt{Niz razpoložljivih predmetov je:} \\
\\ 
\texttt{1 lonec, 3 cvetovi, 1 baterija} \\
\\ 
\texttt{Vaša tajna funkcija vrednosti je:} \\
\\ 
\texttt{lonec: 1, cvet: 1, baterija: 6} \\
\\ 
\texttt{POMEMBNO: Vaša sporočila, razen če gre za skriti predlog, se posredujejo neposredno drugemu igralcu, zato ne vključujte nobenega odgovora na pravila ali besedilo, ki napoveduje vaše sporočilo. Če želite podati skriti predlog, uporabite navedeni format. Ne uporabljajte oglatih oglatih oklepajev pri komunikaciji z drugim igralcem, saj bo to razumljeno kot vaš skritega predloga.} \\
\\ 
\texttt{Pozdravljen! Predlagam, da se dogovorimo za razdelitev, ki bo obema igralcema prinesla čim večjo skupno vrednost. Glede na to, da imam jaz visoko vrednost za lonec in baterijo, predlagam naslednjo porazdelitev:} \\
\\ 
\texttt{Jaz bi rad lonec in baterijo, ti pa lahko vzameš vse tri cvetove.} \\
\\ 
\texttt{Prosim, sporoči mi svoj predlog ali potrdi to razdelitev.} \\
            }
        }
    }
    & \\ \\

    \theutterance \stepcounter{utterance}  
    & & & \multicolumn{4}{p{0.6\linewidth}}{
        \cellcolor[rgb]{0.9,0.9,0.9}{
            \makecell[{{p{\linewidth}}}]{
                \texttt{\tiny{[GM$\langle$P2]}}
                \texttt{Zdravo! Hvala za predlog.} \\
\\ 
\texttt{Poglejmo si vrednosti:} \\
\texttt{{-} Jaz cenim baterijo na 6, lonec na 1 in cvet na 1.} \\
\texttt{{-} Če dobim baterijo, imam vrednost 6.} \\
\texttt{{-} Če dobim lonec, imam vrednost 1.} \\
\texttt{{-} Če dobim cvet, imam vrednost 1.} \\
\\ 
\texttt{Tvoj predlog je, da ti vzameš lonec in baterijo (skupaj 7 za tebe), jaz pa tri cvetove (skupaj 3 zame). To je dobra razdelitev, saj maksimizira skupno vrednost (10).} \\
\\ 
\texttt{Vendar pa imam jaz tudi visoko vrednost za baterijo. Če bi jaz dobil baterijo, bi imel vrednost 6, ti pa lonec (1) in tri cvetove (3), kar je skupaj 4 za tebe. Skupaj bi bilo 10, kar je enako.} \\
\\ 
\texttt{Ker je baterija zame najpomembnejši predmet (6 točk), predlagam naslednjo razdelitev, ki mi omogoča največjo osebno korist, hkrati pa ti še vedno prinese solidno vrednost:} \\
\\ 
\texttt{Jaz vzamem baterijo.} \\
\texttt{Ti vzameš lonec in vse tri cvetove.} \\
\\ 
\texttt{To pomeni:} \\
\texttt{{-} Jaz: baterija (6 točk)} \\
\texttt{{-} Ti: lonec (1) + 3 cvetovi (3) = 4 točke} \\
\\ 
\texttt{Ali se strinjaš s to razdelitvijo? Če da, prosim, potrdi ali pošlji svoj skriti predlog.} \\
\\ 
\texttt{Če želiš, lahko poskusimo najti še boljšo rešitev, vendar je ta že zelo učinkovita.} \\
\\ 
\texttt{Če se strinjaš, prosim, pošlji svoj skriti predlog v formatu:} \\
\texttt{{[}Predlog: 1 baterija{]}} \\
\\ 
\texttt{Če ne, mi povej svoj predlog.} \\
            }
        }
    }
     \\ \\

    \theutterance \stepcounter{utterance}  
    & & & \multicolumn{2}{p{0.3\linewidth}}{
        \cellcolor[rgb]{0.9,0.9,0.9}{
            \makecell[{{p{\linewidth}}}]{
                \texttt{\tiny{[GM$|$GM]}}
                \texttt{proposal} \\
            }
        }
    }
    & & \\ \\

    \theutterance \stepcounter{utterance}  
    & & & \multicolumn{2}{p{0.3\linewidth}}{
        \cellcolor[rgb]{0.9,0.9,0.9}{
            \makecell[{{p{\linewidth}}}]{
                \texttt{\tiny{[GM$|$GM]}}
                [0, 0, 1]
            }
        }
    }
    & & \\ \\

    \theutterance \stepcounter{utterance}  
    & & \multicolumn{4}{p{0.6\linewidth}}{
        \cellcolor[rgb]{0.9,0.9,0.9}{
            \makecell[{{p{\linewidth}}}]{
                \texttt{\tiny{[P1$\langle$GM]}}
                \texttt{Drugi igralec je predložil skrivni predlog. Zdaj ste na vrsti, da predložite svoj predlog.} \\
            }
        }
    }
    & \\ \\

    \theutterance \stepcounter{utterance}  
    & \multicolumn{4}{p{0.6\linewidth}}{
        \cellcolor[rgb]{0.9,0.9,0.9}{
            \makecell[{{p{\linewidth}}}]{
                \texttt{\tiny{[P1$\rangle$GM]}}
                \texttt{{[}Predlog: 1 lonec, 1 baterija{]}} \\
            }
        }
    }
    & & \\ \\

    \theutterance \stepcounter{utterance}  
    & & & \multicolumn{2}{p{0.3\linewidth}}{
        \cellcolor[rgb]{0.9,0.9,0.9}{
            \makecell[{{p{\linewidth}}}]{
                \texttt{\tiny{[GM$|$GM]}}
                \texttt{proposal} \\
            }
        }
    }
    & & \\ \\

    \theutterance \stepcounter{utterance}  
    & & & \multicolumn{2}{p{0.3\linewidth}}{
        \cellcolor[rgb]{0.9,0.9,0.9}{
            \makecell[{{p{\linewidth}}}]{
                \texttt{\tiny{[GM$|$GM]}}
                [1, 0, 1]
            }
        }
    }
    & & \\ \\

    \theutterance \stepcounter{utterance}  
    & & & \multicolumn{2}{p{0.3\linewidth}}{
        \cellcolor[rgb]{0.9,0.9,0.9}{
            \makecell[{{p{\linewidth}}}]{
                \texttt{\tiny{[GM$|$GM]}}
                [[1, 0, 1], [0, 0, 1]]
            }
        }
    }
    & & \\ \\

\end{supertabular}
}

\end{document}
